\section{Key Contributions}

The primary contributions of this thesis are:

\begin{enumerate}
    \item \textbf{Novel Metalinguistic Integration}: We present a novel integration of metalinguistic features (theme and tone) into article-level political bias classification, demonstrating significant performance improvements particularly for challenging center-leaning content.

    \item \textbf{Comprehensive LLM Evaluation}: We conduct an evaluation of state-of-the-art large language models (including GPT-4o, LLaMA-3, and o3-mini) on political bias classification, examining both zero-shot prompting and fine-tuning approaches.

    \item \textbf{Enhanced Dataset and Baseline}: We introduce an expanded version of the AllSides dataset with 47\% additional annotated articles and establish an informal human performance baselines through a small annotation study.

    \item \textbf{Methodological Advances}: We develop improved continued pre-training techniques that achieve new state-of-the-art performance on political bias classification while maintaining robustness to source leakage. Our best model achieves an F1 score of 62.19 on the test set, a 1.59-point improvement over the previous state-of-the-art model POLITICS published by \citet{liu2022politics}.

    \item \textbf{Causal Inference Analysis}: We provide the first causal analysis of theme-tone-ideology relationships in political text classification at both topic and community levels, comparing human expert judgments with automated system outputs across four annotation paradigms. Our community-level analysis reveals that fine-tuned LLMs produce significant causal effects absent in human annotations, evidence of shortcut learning that is invisible to standard F1-based evaluation, and demonstrates mediation analysis as a diagnostic framework for annotation trustworthiness.
\end{enumerate}
